\section{Completeness}
\label{sec:fol-metatheory-03}

\begin{theorem}[Strong completeness of classical first-order logic]
\label{thm:fol-completeness}
Let \(\Lang\) be a signature and assume the Axiom of Choice in the
ambient set theory.  For every
\(\Gamma\cup\set{\varphi}\subseteq\Fm(\Lang)\),
\[
 \Gamma\entails\varphi
 \quad\Longrightarrow\quad
 \Gamma\derives_{\ND_{\FOL}}\varphi.
\]
\end{theorem}
\begin{proof}
We prove the contrapositive.  If \(\Gamma\notderives\varphi\), then
\(\Sigma=\Gamma\cup\set{\neg\varphi}\) is consistent.  For suppose a
derivation \(D\) of \(\bottom\) has open assumptions in \(\Sigma\).  If no
open occurrence of \(\neg\varphi\) appeared in \(D\), then \(D\) would
witness \(\Gamma\derives\bottom\), and \(\bottom E\) would give
\(\Gamma\derives\varphi\); otherwise mark all open occurrences of
\(\neg\varphi\) in \(D\) with a common fresh label and close \(D\) by
\(\mathsf{RAA}\) at the root, discharging every such occurrence, to
derive \(\varphi\) from \(\Gamma\).  Either way
\(\Gamma\derives\varphi\), contrary to supposition.  Apply the
Henkin--Lindenbaum lemma to obtain a maximal consistent witness set
\(\Delta\supseteq\Sigma\).  Its term structure and canonical assignment satisfy
every member of \(\Delta\) by the truth lemma.  In particular,
\(\Meta_\Delta,s_0\entails\Gamma\) and
\(\Meta_\Delta,s_0\entails\neg\varphi\).  Thus
the reduct \(\Meta_\Delta\restr\Lang\), with the same assignment \(s_0\),
satisfies \(\Gamma\) and falsifies \(\varphi\) by
\cref{lem:reduct-invariance}.  It is a \(\Lang\)-structure, so
\(\Gamma\notentails\varphi\) under the semantic consequence relation stated
in the theorem.
\end{proof}

Together with soundness, completeness gives
\[
 \Gamma\derives_{\ND_{\FOL}}\varphi
 \quad\text{iff}\quad
 \Gamma\entails\varphi.
\]
This equivalence concerns the specified calculus and the set-based semantics
with nonempty domains.  It does not identify syntactic derivability
with semantic truth as objects.

\begin{corollary}[Completeness for countable signatures, without choice]
\label{cor:fol-completeness-countable}
Let \(\Lang\) be a signature whose set of function and relation symbols
is countable in the sense of \cref{sec:set-prerequisites-05}, finite
signatures included.  Then, for every
\(\Gamma\cup\set{\varphi}\subseteq\Fm(\Lang)\),
\[
 \Gamma\entails\varphi
 \quad\Longrightarrow\quad
 \Gamma\derives_{\ND_{\FOL}}\varphi,
\]
and no use is made of the Axiom of Choice.  This is the
countable-language version of the completeness theorem promised in
\cref{sec:set-prerequisites-05}.
\textit{Status: complete proof.}
\end{corollary}
\begin{proof}
Run the proof of \cref{thm:fol-completeness} unchanged, replacing each
input that the Axiom of Choice supplied with an explicit one.  Choice
enters that proof only through \cref{lem:henkin-lindenbaum}, at exactly
two points: the cardinal bookkeeping at
\(\kappa=\max(\abs{\Lang},\aleph_0)\) through
\cref{cor:cardinal-assignment}, \cref{thm:infinite-cardinal-finite-coding},
and \cref{cor:syntax-counting}; and the bijective enumeration of
\(\Fm(\Lang^H)\) handed to \cref{lem:transfinite-lindenbaum}, whose own
construction is choice-free given the enumeration, as its statement
records.  The enumeration of the existential formulas
\(\langle E_\alpha\rangle\) is an instance of the same second point.

Both inputs are explicit here.  Fix an injection \(\iota\) of the symbol
set of \(\Lang\) into \(\N\); one exists by the definition of countable
set, and fixing a single witness of one existential statement is not an
application of the Axiom of Choice.  Take \(\kappa=\aleph_0\) and add the
witness constants \(\langle d_k:k\in\N\rangle\).  The alphabet of
\(\Lang^H\) is the union of three explicitly enumerated families: the
finitely many logical symbols and punctuation marks together with the
variables in the fixed order \(x_0,x_1,\ldots\) declared in
\cref{ch:fol-syntax-semantics}; the symbols of \(\Lang\) in the order
given by \(\iota\); and the constants \(d_k\) in index order.  Sending
the \(k\)th member of the three families to \(3k\), \(3k+1\), and
\(3k+2\) respectively injects the alphabet into \(\N\), and the string
coding of \cref{ch:strings}---iterating \(\operatorname{pair}_{\N}\) as
in the countability-of-finite-strings lemma---then assigns every formula
of \(\Lang^H\) a distinct natural number code.  Define \(\varphi_n\) to
be the formula whose code is the \(n\)th smallest member of the set of
formula codes; recursion on \(\N\) with the least-number principle
produces this bijective enumeration of \(\Fm(\Lang^H)\) without choice,
and its restriction to the existential formulas enumerates them likewise.
The transfinite recursions of \cref{lem:henkin-lindenbaum} and
\cref{lem:transfinite-lindenbaum} then run along \(\omega\) and
\(\omega+1\) with functional clauses.  At \(\kappa=\aleph_0\) the
cardinal-arithmetic appeals reduce to finite counting and to the displayed
coding: before any stage \(n\) of the witness recursion, only finitely
many witness implications have been added, so they contain only finitely
many of the \(d_k\) and a least available one exists; and the code
injection directly witnesses that \(\Fm(\Lang^H)\) is countably
infinite.

No other ingredient of the completeness proof invokes a choice principle:
\cref{lem:new-constant-elimination} and
\cref{lem:witness-extension} manipulate single finite derivations; the
term structure interprets symbols on equivalence classes with
representative independence proved by
\cref{lem:equality-congruence}; the canonical assignment is
\(s_0(x)=[x]\); and the truth lemma and
\cref{lem:reduct-invariance} are inductions all of whose clauses are
explicit.  The
contrapositive argument of \cref{thm:fol-completeness} therefore goes
through verbatim for countable \(\Lang\) with no appeal to the Axiom of
Choice.
\end{proof}

\begin{theorem}[Completeness of the first-order sequent calculus]
\label{thm:fol-lk-completeness}
Assume the Axiom of Choice in the ambient set theory.  For all finite
multisets \(\Gamma,\Delta\) of \(\Lang\)-formulas, the sequent
\(\Gamma\Rightarrow\Delta\) is derivable in \(\LK_{\FOL}\) iff for every
\(\Lang\)-structure \(\Meta\) and assignment \(s\): if
\(\Meta,s\entails\gamma\) for every
\(\gamma\in\operatorname{supp}(\Gamma)\), then
\(\Meta,s\entails\delta\) for some
\(\delta\in\operatorname{supp}(\Delta)\).
\textit{Status: proof sketch.}
\end{theorem}
\begin{proof}[Proof sketch]
The exact scope is the cut-admitting calculus of
\cref{ch:fol-deduction}.  Left to right is \cref{thm:fol-lk-soundness},
whose satisfaction condition on a sequent is exactly the one displayed
here; at first order the semantic data are a structure together with an
assignment.

For right to left, first record the translation subroutine.  Induction on
a \(\ND_{\FOL}\)-derivation translates the
propositional rules as in the propositional completeness proof, translates
\(\forall I\) and \(\exists E\) using the corresponding eigenvariable rules,
translates \(\forall E\), \(\exists I\), and equality rules directly, and
translates the \(\alpha\) rule by an induction on the generating
alpha-equivalence steps using the sequent quantifier rules.  Cut composes
translated elimination steps.  Bound-variable alpha-equivalence and the
separate eigenvariable freshness conditions are preserved by the relevant
subinductions.  The induction establishes that if
\(\Gamma_0\derives_{\ND_{\FOL}}\varphi\) for a finite set \(\Gamma_0\),
then the sequent \(\Gamma_0\Rightarrow\varphi\) is derivable in
\(\LK_{\FOL}\).  The individual translation trees and those subinductions
are omitted.

Now suppose the semantic condition holds, and list
\(\operatorname{supp}(\Delta)=\set{\delta_1,\ldots,\delta_m}\) with
\(m\geq0\).
\begin{enumerate}
  \item Let \(S\) be the set of members of
        \(\operatorname{supp}(\Gamma)\) together with
        \(\neg\delta_1,\ldots,\neg\delta_m\).  A structure and
        assignment satisfying all of \(S\) would satisfy every member
        of \(\operatorname{supp}(\Gamma)\), hence some \(\delta_j\)
        together with \(\neg\delta_j\).  So no structure and
        assignment satisfies \(S\), and \(S\entails\bottom\).
  \item \Cref{thm:fol-completeness} gives
        \(S\derives_{\ND_{\FOL}}\bottom\).  The open assumptions of
        that derivation form a finite subset of \(S\), say
        \(\gamma_1,\ldots,\gamma_k\in\operatorname{supp}(\Gamma)\) and
        \(\neg\delta_{j_1},\ldots,\neg\delta_{j_r}\).  The translation
        subroutine turns it into an \(\LK_{\FOL}\)-derivation of
        \(\gamma_1,\ldots,\gamma_k,
        \neg\delta_{j_1},\ldots,\neg\delta_{j_r}\Rightarrow\bottom\).
  \item Cut with the initial \(\bottom\)-sequent
        \(\bottom\Rightarrow\), the instance of
        \(\Gamma,\bottom\Rightarrow\Delta\) with both contexts empty,
        empties the succedent.
  \item For each \(j\), the sequent
        \(\Rightarrow\delta_j,\neg\delta_j\) is derivable: start from
        the identity sequent \(\delta_j\Rightarrow\delta_j\) and apply
        \(\neg R\).  Cutting it against the current sequent, with cut
        formula \(\neg\delta_j\),
        \[
         \Rule[Cut]{\Rightarrow\delta_j,\neg\delta_j\qquad
          \Gamma',\neg\delta_j\Rightarrow\Pi}
          {\Gamma'\Rightarrow\delta_j,\Pi},
        \]
        moves \(\delta_j\) into the succedent; after \(r\) such cuts
        the derived sequent is
        \(\gamma_1,\ldots,\gamma_k\Rightarrow
        \delta_{j_1},\ldots,\delta_{j_r}\).
  \item Weakening (\(W_L\), \(W_R\)) adds the unused members of
        \(\Gamma\) and \(\Delta\), and contraction (\(C_L\),
        \(C_R\)) merges duplicates, giving exactly the multiset
        sequent \(\Gamma\Rightarrow\Delta\).
\end{enumerate}
When \(m=0\), the hypothesis asserts that no structure and assignment
satisfies every member of \(\operatorname{supp}(\Gamma)\), step 4
performs no cuts, and the same route derives
\(\Gamma\Rightarrow\Delta\) with empty succedent; an empty \(\Gamma\) is
likewise covered.  The route uses cut, which is primitive in the
displayed calculus, and makes no cut-free claim; extending the
subroutine's translation induction to multi-succedent targets is
unnecessary, because the reduction reuses it verbatim.  The status is a
sketch because the individual translation trees and subinductions of the
subroutine remain omitted; the argument establishes completeness but not
cut elimination.
\end{proof}

